\documentclass[12pt]{article}
\usepackage[utf8]{inputenc}
\usepackage[T1]{fontenc}
\usepackage[margin=1in]{geometry}
\usepackage{setspace}
\usepackage{xcolor}
\usepackage[colorlinks=true, urlcolor=blue!60!black, linkcolor=blue!60!black]{hyperref}
\setlength{\parindent}{2em}
\setlength{\parskip}{1.25em}
\renewcommand{\baselinestretch}{1.0}
\begin{document}
\begin{center}
{\LARGE \textit{Duncan v.\ Louisiana}, 391 U.S.\ 145 (1968)\par}
\vspace{0.8em}
{\large Case Brief --- Four-Page Edition\par}
\end{center}
\vspace{1em}

\section{Memory Jogger}

The Court incorporated the Sixth Amendment right to a criminal jury trial against the states through the Fourteenth Amendment's Due Process Clause, holding that any state prosecution for a serious offense, including a misdemeanor punishable by up to two years' imprisonment, must afford a jury (\href{https://supreme.justia.com/cases/federal/us/391/145/}{391 U.S. 145}).

\section{Detailed Case Facts}

On October 18, 1966, Gary Duncan, a nineteen-year-old Black man, stopped his car in Plaquemines Parish, Louisiana, to urge two younger cousins, who had recently transferred to a formerly all-white school, to leave an encounter with four white boys. The white witnesses testified that Duncan slapped one boy, Herman Landry, on the elbow; Duncan and his cousins said he merely touched him.

Duncan was charged with simple battery, a misdemeanor punishable by up to two years' imprisonment and a \$300 fine. The Louisiana Constitution granted jury trials only where capital punishment or imprisonment at hard labor could be imposed, so the judge denied Duncan's request for a jury, convicted him at a bench trial, and sentenced him to 60 days and a \$150 fine. The prosecution unfolded amid school desegregation in a parish dominated by segregationist Leander Perez; civil rights lawyers including Richard Sobol and Anthony Amsterdam represented Duncan.

\section{Procedural History}

The Louisiana Supreme Court denied review, finding no error of law. Duncan appealed to the U.S. Supreme Court, which heard the case alongside \textit{Bloom v.\ Illinois} and on May 20, 1968, reversed and remanded.

\section{Judicial Votes}

Reversed and remanded, \textbf{7--2}.
\begin{itemize}
\item \textbf{Majority:} Justice White, joined by Chief Justice Warren and Justices Black, Douglas, Brennan, Fortas, and Marshall.
\item \textbf{Concurring:} Justice Black, joined by Justice Douglas; Justice Fortas (in an opinion covering \textit{Duncan} and \textit{Bloom}).
\item \textbf{Dissenting:} Justice Harlan, joined by Justice Stewart.
\end{itemize}

\section{Holding}

Because trial by jury in criminal cases is fundamental to the American scheme of justice, due process guarantees a jury in every state prosecution that would fall within the Sixth Amendment if tried federally. Petty offenses remain exempt, but the maximum authorized penalty, not the sentence imposed, is the chief measure of seriousness, and a crime punishable by two years is serious. The Court did not fix the exact petty-offense line (later set at six months in \textit{Baldwin v.\ New York}), and Footnote 30's remarks on jury size and unanimity were dicta.

\section{Analysis of Opinions}

\textbf{Majority (White).} White reframed incorporation. Rather than asking, as \textit{Palko} had, whether any fair system could exist without the safeguard, the Court asked whether it is fundamental in the American system, which is built around the common-law jury. He traced the jury from Magna Carta and Blackstone through the Declaration of Independence, Article III, the Sixth Amendment, and every state constitution. Functionally, the jury guards against government oppression by requiring community participation before liberty is taken. He set aside contrary language in \textit{Maxwell v.\ Dow}, \textit{Snyder}, and \textit{Palko}, and cited Kalven and Zeisel's \textit{The American Jury} on jurors' competence. On seriousness, he treated the statutory maximum as the objective gauge, noting that 49 states limited nonjury trials to crimes punishable by a year or less.

\textbf{Black (concurring).} Reaffirming his \textit{Adamson} dissent, Black argued that the Fourteenth Amendment's framers, Bingham and Howard, meant to apply the whole Bill of Rights to the states, rebutted Charles Fairman's critique, and preferred selective incorporation to a fairness test that gives judges unconfined power.

\textbf{Fortas (concurring).} Fortas joined the result but rejected jot-for-jot incorporation: the core jury right need not carry every federal rule, such as twelve jurors or unanimity, because the Constitution creates a federal union, not a monolith.

\textbf{Harlan (dissenting).} Selective incorporation is an illogical compromise. Relying on Fairman, Harlan read due process as requiring fundamental fairness, not the Bill of Rights jot-for-jot. Bench trials are not inherently unfair, as jury waivers and English summary trials show, and uniform federal rules would either bind the states in a ``constitutional straitjacket'' or water down federal rights. States should remain laboratories.

\section{Examples: Future Applications}

\begin{enumerate}
\item \textbf{Same-Side (Jury Required): One-Year Shoplifting.} A state punishes shoplifting by up to twelve months in jail but denies a jury. The authorized maximum, not a promised 30-day sentence, controls, so a jury is required.
\item \textbf{Same-Side (Jury Required): Uncapped Contempt.} A judge sentences a lawyer to eighteen months for criminal contempt under a rule with no statutory maximum. Under \textit{Bloom}, the sentence imposed measures seriousness, so a jury is required.
\item \textbf{Opposite-Side (No Jury): Reckless Driving.} An offense carrying at most 90 days and a \$250 fine falls within the preserved petty-offense exception.
\item \textbf{Opposite-Side (No Jury): Civil Penalties.} A \$50,000 civil forfeiture action with no incarceration is not a criminal case, and the Seventh Amendment remains unincorporated.
\item \textbf{Fence-Sitter: Stacked Petty Charges.} Four 90-day disorderly-conduct counts create 360 days of exposure. \textit{Duncan}'s concern with liberty suggests a jury, but its focus on the penalty for a particular crime points the other way, as \textit{Lewis v.\ United States} later held.
\end{enumerate}

\section{Critique}

\textbf{Doctrinal and progressive critiques.} Harlan exposed a tension in selective incorporation: no textual explanation shows why the criminal jury is in while grand juries and civil juries are out. His warning about watering down proved prescient when \textit{Williams v.\ Florida} and \textit{Apodaca v.\ Oregon} diluted jury size and unanimity, until \textit{Ramos v.\ Louisiana} restored unanimity in 2020. Civil rights scholars, and Gary Duncan himself, note that the jury right did not guarantee representative juries, which awaited \textit{Taylor v.\ Louisiana}; that Duncan escaped retrial only through a federal injunction; and that plea bargaining and \textit{McKeiver} limited the decision's reach.

\textbf{Originalist critiques.} Many originalists accept incorporation of the jury right but argue that the Privileges or Immunities Clause, not due process, was the proper vehicle, as Justice Thomas argued in \textit{McDonald} and \textit{Ramos}, and they fault \textit{Williams} for abandoning the twelve-person jury. Federalism-minded conservatives echo Harlan's concern for state experimentation.

\section{Key Quotations}

\begin{quote}
``Because we believe that trial by jury in criminal cases is fundamental to the American scheme of justice, we hold that the Fourteenth Amendment guarantees a right of jury trial in all criminal cases which --- were they to be tried in a federal court --- would come within the Sixth Amendment's guarantee.''
\end{quote}
\noindent --- Justice White, Opinion of the Court, 391 U.S. at 149.

\begin{quote}
``Providing an accused with the right to be tried by a jury of his peers gave him an inestimable safeguard against the corrupt or overzealous prosecutor and against the compliant, biased, or eccentric judge. If the defendant preferred the common-sense judgment of a jury to the more tutored but perhaps less sympathetic reaction of the single judge, he was to have it.''
\end{quote}
\noindent --- Justice White, Opinion of the Court, 391 U.S. at 156.

\begin{quote}
``The question thus is whether given this kind of system a particular procedure is fundamental --- whether, that is, a procedure is necessary to an Anglo-American regime of ordered liberty.''
\end{quote}
\noindent --- Justice White, Opinion of the Court, 391 U.S. at 149--150 n.14.

\end{document}
